\chapter{Future Work}\label{ch:future}

Future work is deliberately organised into \emph{three parallel tracks}
rather than a single laundry list.  The tracks are complementary:
Track~1 winds the pipeline's human-in-the-loop dependency \emph{down},
while Tracks~2 and 3 widen its coverage \emph{outward} to new languages
and new social-media platforms.  Every deliverable is admissible only
when the previous milestone in its track is measurably successful; the
guiding principle in all three tracks is that human effort produces the
training signal that unlocks the next step.

\begin{center}
\small
\begin{tabularx}{\linewidth}{@{}p{2.4cm} p{2.6cm} X@{}}
\toprule
\textbf{Track} & \textbf{Direction} & \textbf{Question it answers} \\
\midrule
Track 1 & HITL reduction (deeper) &
How do we go from today's 100\,\% analyst-supervised replies to a
mostly-autonomous replier, without a silent quality drop? \\
Track 2 & Bilingual $\to$ multilingual (wider) &
How do we cover Indian retail communities that today's English-only
filter silently drops? \\
Track 3 & Multi-platform (wider) &
How do we extend beyond Reddit to X, Facebook, Instagram, and
app-store reviews so brand health reflects the whole social surface? \\
\bottomrule
\end{tabularx}
\end{center}

\section{Track 1 --- HITL Reduction (Four-Horizon Roadmap)}\label{sec:future:hitl}

Track~1 is the deepest of the three because it changes the pipeline's
operating mode over time.  It is expressed as four horizons of
declining HITL involvement, each shipping a concrete artefact that
becomes the training signal for the next horizon.

\definecolor{hzGreen}{HTML}{2E7D32}
\definecolor{hzBlue}{HTML}{0071DC}
\definecolor{hzPurple}{HTML}{6A1B9A}
\definecolor{hzAmber}{HTML}{E65100}

\renewcommand{\arraystretch}{1.35}
\begin{table}[htbp]
\centering
\small
\begin{tabularx}{\linewidth}{@{}p{2.6cm} p{2.1cm} X@{}}
\toprule
\textbf{Horizon} & \textbf{HITL} & \textbf{Deliverable} \\
\midrule
\rowcolor{zebra}
\textcolor{hzGreen}{\textbf{Now $\to$ 1 month}} &
100\,\% &
\textbf{Grow the gold set.}  Grow the labelled gold set to
$\sim$500 posts using the HITL feedback the team is already producing;
rerun 5-fold OOF to check the F1 delta. \\

\textcolor{hzBlue}{\textbf{1 $\to$ 3 months}} &
100\,\% &
\textbf{Automate the retrain loop.}  Wire the monthly ModernBERT
retrain into a Cron / Airflow job with an OOF-F1 promotion gate --- new
checkpoints ship only when they beat the current one on held-out folds. \\

\rowcolor{zebra}
\textcolor{hzPurple}{\textbf{3 $\to$ 6 months}} &
$\sim$80\,\% &
\textbf{Walmart-trained model + shared DB.}  Fine-tune Mistral 7B on
our own posted-reply corpus $\to$ a Walmart-trained replier that drops
the paid GPT-4o dependency, runs fully in-house, and can be self-hosted.
Move the feedback table to a shared managed database so multiple
analysts can work concurrently. \\

\textcolor{hzAmber}{\textbf{6 $\to$ 12 months}} &
$\sim$30\,\% $\to$ 0\,\% &
\textbf{Confidence-gated auto-reply.}  Introduce a confidence gate on
the reply generator: top-confidence drafts are queued for auto-send and
the analyst only reviews the mid- and low-confidence tail.  As
agreement stays high, gradually raise the auto-send threshold --- HITL
winds down to sampling / audit only. \\
\bottomrule
\end{tabularx}
\caption{Track 1 --- Four-horizon HITL reduction roadmap.  Each row
lists the target HITL involvement and the deliverable that unlocks the
next horizon.}\label{tab:future:horizons}
\end{table}
\renewcommand{\arraystretch}{1.0}

\paragraph{Guiding principle.}
\emph{HITL is the ladder, not the ceiling} --- every stage produces the
training signal that lets the next stage automate more.  A horizon is
declared complete only when its OOF-F1 (for classification) or
analyst-edit-distance (for replies) meets the promotion gate on a
statistically meaningful sample; otherwise the horizon is extended
rather than skipped.

\section{Track 2 --- Bilingual (Hindi + English) $\to$ Multilingual}\label{sec:future:multilingual}

\paragraph{Why now.}
Indian retail-adjacent communities on Reddit --- \texttt{r/india},
\texttt{r/IndiaOffline}, \texttt{r/IndianRealEstate},
\texttt{r/DebateIndianEconomy}, plus the shopping-heavy
sub-communities of \texttt{r/personalfinanceindia} --- account for an
estimated 30\,\% of Walmart-family mentions (Walmart owns Flipkart,
PhonePe, and Myntra in India).  Today's ingestion layer runs a
\texttt{langdetect} filter that silently drops anything not classified
as English.  Any post in Hindi, Tamil, Telugu, Bengali, or
code-mixed \emph{Hinglish} is invisible to the rest of the pipeline ---
including brand health, alerts, and the reply queue.  This is a
measurement failure, not just a coverage gap: our current
``Walmart brand sentiment'' score is silently biased toward
English-speaking, US-centric communities.

\paragraph{Concrete deliverables.}
\begin{itemize}
\item \textbf{Detector.}  Extend \texttt{langdetect} to accept
      \texttt{\{hi, en\}} first, then add
      \texttt{\{ta, te, bn\}} once the analyst pool for each
      language is in place.  Handle code-mixed Hinglish (dominant on
      Indian Reddit) with a mixed-script tokenizer.
\item \textbf{Two viable modelling paths}, decided empirically:
      \emph{(a)} translate posts to English at ingestion time via
      NLLB-200 (Meta, open-weights, 3.3\,B) so downstream sentiment
      and aspect models stay English-only, \emph{or} \emph{(b)}
      fine-tune a second ModernBERT checkpoint on native-script data
      and route by detected language.  Path (a) is cheaper to ship;
      path (b) is more accurate for low-resource languages.  Pilot
      both on a $\sim$200-post Hindi gold set before committing.
\item \textbf{Taxonomy extension.}  The seven customer aspects and
      five workforce aspects were validated on English Reddit only.
      Indic retail idioms (``bill karo'' for billing, ``delivery
      kab'' for late shipping, ``paisa wapas'' for refund) need to
      be added or the aspect NLI head will miss them.
\item \textbf{Native-speaker HITL loop.}  Reuse the existing Review \&
      Validate UI unchanged; only the analyst pool changes.  The
      warm-loop retrain in Section~\ref{sec:future:hitl} then covers
      Hindi automatically once the gold set has $\sim$100 Hindi posts.
\end{itemize}

\paragraph{How this track ties into Track 1.}
Track~2 rides on top of Track~1's horizon 1--3 months (automate
the retrain loop) and horizon 3--6 months (Walmart-trained model).
Once the monthly retrain is automated, adding Hindi is a matter of
appending Hindi gold-set posts to the training pool; once the
Walmart-trained Mistral checkpoint exists, its multilingual variant
(Mistral supports 15+ languages natively) becomes the natural
production reply generator.

\section{Track 3 --- Multi-Platform (Reddit $\to$ X / Facebook / Instagram / app-stores)}\label{sec:future:multiplatform}

\paragraph{Why now.}
Reddit alone answers \emph{what deliberate, long-form customers say}.
It misses:
\begin{itemize}
\item \textbf{X (formerly Twitter)} --- real-time viral complaints;
      brand-tagged tweets are the fastest signal of an outage or
      recall.
\item \textbf{Instagram} --- image-first customer content
      (unboxing videos, product-defect photos); this is where our
      four-pass vision pipeline (Section~\ref{sec:vision}) actually
      shines because the caption / OCR / merge stack was built for
      exactly this format.
\item \textbf{Facebook} --- the age-demographic complement to
      Reddit; older, family-focused, and often the first channel
      customers reach for on returns or in-store complaints.
\item \textbf{Google Play / Apple App Store reviews} --- the
      highest-intent signal in the entire retail social surface,
      because every review carries a native
      \texttt{star\_rating} field (1--5) that is effectively free
      ground truth for sentiment.
\end{itemize}
A brand-health score built on Reddit-only is systematically biased
toward what one deliberate, English-speaking subset of the internet
says.

\paragraph{Concrete deliverables.}
\begin{itemize}
\item \textbf{Ingestion adapters}, one per platform, each behind a
      common \texttt{RawPost} interface so the analysis pipeline
      downstream is platform-agnostic: X API v2, Meta Graph API
      (Facebook + Instagram), Google Play scraping via
      \texttt{google-play-scraper}, App Store via the RSS feed.
\item \textbf{Per-platform pre-processing.}  X posts need retweet-
      lineage dedup (a retweet is not a new opinion); Instagram
      posts route through the vision pipeline first (caption + image
      OCR) before hitting text sentiment; app-store reviews carry a
      native \texttt{star\_rating} that is stored alongside our
      ModernBERT sentiment for direct calibration checks.
\item \textbf{Extended priority classifier.}  The existing P1/P2 rules
      (Section~\ref{sec:alerts}) need platform-specific extensions:
      an X post with $\geq$~100 retweets is auto P1; an app-store
      1-star surge (rate-of-change on the daily 1-star count) is P1;
      an Instagram brand-tagged post from a $\geq$~10\,k-follower
      account is P1.  Reddit rules are unchanged.
\item \textbf{Unified schema.}  The
      \texttt{raw\_posts.source\_platform} column is already stubbed
      in the current schema (it takes the value \texttt{"reddit"}
      for every row today).  Populating it correctly is a one-line
      change per adapter and lets Brand Health / Insights aggregate
      across platforms out of the box.
\end{itemize}

\paragraph{Known risks.}
X's post-2023 API tiering makes free ingestion impractical (Basic
tier is \$100/month for $\sim$10\,k posts/month, which is enough for a
pilot but not for production).  The Meta Graph API requires
page-level admin permissions --- workable inside Walmart but a
blocker for open replication.  Instagram image volumes will
overwhelm the current one-image-at-a-time vision pass and will need
batched inference on a GPU pool.

\section{Beyond the Three Tracks}\label{sec:future:beyond}

The following items are on the wishlist but were deliberately excluded
from the three tracks because they lack a clear next step grounded in
the current pipeline.  They are listed here so reviewers can see the
design space we chose \emph{not} to commit to.

\subsection*{Model-level extensions}
\begin{itemize}
\item Train a joint sentiment + aspect head with a shared encoder to
      reduce inference cost by a further $\sim$30\,\%.
\item Distil ModernBERT into a $\sim$50\,M-parameter student for
      CPU-only inference at the edge.
\item Add reasoning-augmented VLM captioning (LLaVA-Next 1.6\,B or
      SmolVLM) to the vision layer for step-through descriptions of
      complex screenshots.
\end{itemize}

\subsection*{Product-level extensions}
\begin{itemize}
\item Predictive P1 forecast: seasonal-decomposition on daily P1
      counts to forecast alert volume 24--48 hours ahead.
\item Sentiment-drift alerts on the Walmart-owned brand set
      (Flipkart, PhonePe, Myntra) as an early warning for
      cross-brand contagion.
\end{itemize}

\subsection*{Operational extensions}
\begin{itemize}
\item Deploy the pipeline as a scheduled Kubernetes CronJob backed by
      the shared managed database introduced in Track 1 horizon 3.
\item Integrate with Walmart's internal ticketing so P1 alerts open
      cases directly.
\end{itemize}
